\section*{3. 对偶簇}

\subsubsection*{3.1.}
设 $I$ 为 $\mathbb P^r$ 中定义 $X$ 的理想. 以 $N=I/I^2$ 表示余法层, 它在 $X$ 上局部自由秩为 $r-n$; 设 $\mathbb P(N)$ 为 $X$ 上由 $\check N$ 中的直线构成的射影丛.(当 $X=\mathbb P^r$ 时, 丛 $\mathbb P(N)$ 为空.)\EditorialFootnote{记号核对: 法文原书第 217 页和英文底本此处把 $N$ 称为余法层, 但 5.1 的正合列和 5.5 的 Chern 类计算改用 $\check N$ 表示余法层、$N$ 表示法层. 这是底本前后不一致的记号, 此处保留而不静默统一. 下文次数计算应按 5.1 与 5.5 的约定理解.} 此丛 $\mathbb P(N)$ 经映射
\[
\nu:\mathbb P(N)\longrightarrow X\times\check{\mathbb P}^{r}.
\]
可等同于 $X\times\check{\mathbb P}^{r}$ 中由满足 $H$ 在 $x$ 处与 $X$ 相切的点对 $(x,H)$ 构成的子簇. 更明确地,
\begin{equation}\tag{3.1.1}
\mathbb P(I/I^2)_x\ni F\longmapsto
\left(x,\text{ 由下列方程定义的超平面 }
\sum_i\frac{\partial F}{\partial X_i}(x)T_i=0\right),
\end{equation}
其中 $F$ 表示 $I/I^2$ 的局部截面. 定义
\begin{equation}\tag{3.1.2}
\varphi:\mathbb P(N)\longrightarrow\check{\mathbb P}^{r}
\end{equation}
为复合
\[
\mathbb P(N)\xrightarrow{\nu}
X\times\check{\mathbb P}^{r}
\xrightarrow{\operatorname{pr}_2}\check{\mathbb P}^{r}.
\]

\subsubsection*{3.1.3.}
以 $\check X$ 表示 $\varphi$ 的像. 它是与 $X$ 不横截相交的超平面所成的集合, 称为 $X$ 关于嵌入 $i:X\hookrightarrow\mathbb P^r$ 的\underline{对偶簇}.

\subsubsection*{命题 3.1.4.}
\underline{对偶簇 $\check X$ 不可约且固有, 维数至多为 $r-1$.}\EditorialFootnote{空集情形: 证明已说明, 若 $X=\mathbb P^r$, 则 $\check X$ 为空. 按通常不将空集称为不可约集的约定, 本命题的不可约性断言应限于 $\check X$ 非空的情形.}

\noindent\underline{证明.}
它是该射影丛在 $\varphi$ 下的像, 当 $X\ne\mathbb P^r$ 时 $\mathbb P(N)$ 本身不可约且维数为 $r-1$, 当 $X=\mathbb P^r$ 时为空.

\subsubsection*{注 3.1.5.}
对 $\mathbb P(N)$ 的解释可以更加精确如下. 设 $\mathcal Y\subset X\times\check{\mathbb P}^{r}$ 为关联子概形, 其在每点 $H\in\check{\mathbb P}^{r}(k)$ 上的纤维为 $X\cap H$. 概形 $\mathcal Y$ 显然光滑, 维数为 $n+r-1$, 因为它是 $X$ 上相对维数 $r-1$ 的射影丛. 因而 $\mathcal Y\to\check{\mathbb P}^{r}$ 使 $\mathcal Y$ 成为虚相对维数 $(n+r-1)-r=n-1$ 的相对完全交.

于是 \underline{Jacobian 子概形}
\[
J_{n-1}(\mathcal Y/\check{\mathbb P}^{r})
\]
(VI, 5.2)恰为 $\mathbb P(N)$. 相对完全交的 Jacobian 子概形的形成与任意基变换相容. 因此对每条直线 $D\subset\check{\mathbb P}^{r}$, Jacobian 子概形 $J_{n-1}(\widetilde X/D)$ 恰为 $\mathbb P(N)\times_{\check{\mathbb P}^{r}}D$, 其中
\[
\widetilde X=\mathcal Y\times_{\check{\mathbb P}^{r}}D,
\]